\label{sec:constant-dilation}
The direct quartic and the compact construction in this section have different
coefficient architectures. Here the coefficients are the complete retained
BFS943/radix67 tuple $(\beta,K_0,M_C,M_F)$, whose largest member has
$592884$ bits. The new equation has eleven strictly positive integer
unknowns, exact degree $174$, and $104$ live binary arithmetic operations.
It retains the native/compiler/refutation calibration. Its full ordinary
proof and all coefficient bytes are included in the companion directory.
No equality with the direct quartic is asserted.

The inspected source-115 comparator has eleven positive unknowns, degree
$211$, and $107$ operations. Its four coefficient bit lengths agree but its
whole-tuple $K_0$ differs. Thus the comparison concerns domain and size,
not an identity between the specialized polynomials. The source-172 equation
with ten positive unknowns, degree $214$, and $87$ operations uses the same
complete tuple as the present construction and remains a distinct tradeoff.

\subsection{The complete polynomial}
All symbols in this section are local. The independent inputs are
\[
 J,F,\alpha,i,T,w,Y,n,\eta,\zeta,\kappa\in\mathbb Z_{>0}.
\]
The circuit JSON retains the historical input name \texttt{s} for $Y$;
this is a renaming of one input, not an additional variable.
Every other symbol below is a polynomial abbreviation. Set $\lambda=210$ and
\begin{align*}
 h&=\beta J,& q&=h+1,\\
 C&=q-F-\alpha,&X&=qw,\\
 k&=\eta+\zeta,&c&=kY+\eta,\\
 A&=\lambda Y(X+1)+1,&\Delta&=A^2-1,\\
 H&=4A-5,&v_0&=XY^2,\\
 v&=\lambda^2v_0,&\ell&=vk,\\
 E&=1+\ell(\ell+k),&B&=(K_0+w)C-F.
\end{align*}
\[
 R=[h(q-F)+\alpha](q^2-1)+(M_C+qM_F)J.
\]
\begin{align*}
 U&=ic^2,&a&=\Delta U^2E,\\
 F_0&=1+a,&G&=EF_0,\\
 f&=1+2a,&Q&=4\Delta aF_0,\\
 b&=cT-R,&V&=fb-c,\\
 L&=bV,&K&=X(3c-2X)+2.
\end{align*}
Define the fixed integer polynomial
\begin{align*}
 \mathcal T(k,x)={}&k-7350xk(k^2-4)\\
 &+5402250x^2k(k^2-4)(27k^2+8)\\
 &-5672362500x^3k(k^2-4)(675k^4-324k^2+160).
\end{align*}
The final registers are
\begin{align*}
 e_*&=\mathcal T(k,v_0)-R-1,&D&=q^7,&M&=hD,\\
 Z&=M(n+\kappa-LB)-DB-h(e_*-q^4Y),\\
 P_0&=Q(Z^2-M^2GV^2)+(M^2G-Z^2),&N&=K-\kappa H.
\end{align*}
The equation is
\[
 \boxed{\quad P=(P_0+1)(3N+1)-1=0.\quad}
\]
There are no extra quantified registers, radicals, variable exponents or
division gates. Since $P_0,N$ are integers, $P=0$ forces both displayed
factors to be units. The second factor is $1$ modulo $3$, and therefore
cannot be $-1$. Consequently
\[
 P=0\quad\Longleftrightarrow\quad P_0=N=0.
\]
The independently integral projector quotient $\kappa$ is essential.

\subsection{An inverse recurrence with integral coefficients}
Let $\psi_b(0)=0$, $\psi_b(1)=1$ and
$\psi_b(m+2)=2b\psi_b(m+1)-\psi_b(m)$. Direct recurrence induction gives
\[
 2\psi_{1+2v}(m)
 =2\sum_{j=0}^{m-1}\binom{m+j}{2j+1}(4v)^j.
\]
Put
\begin{align*}
 C_1(k)&=-\frac{k(k^2-4)}6,\\
 C_2(k)&=\frac{k(k^2-4)(27k^2+8)}{360},\\
 C_3(k)&=-\frac{k(k^2-4)(675k^4-324k^2+160)}{15120}.
\end{align*}
These are the first three corrections in the formal inverse
$k+C_1(k)v+C_2(k)v^2+C_3(k)v^3$. With $\lambda=210$, their denominators
clear before any congruence argument:
\[
 \lambda^2/6=7350,\qquad
 \lambda^4/360=5402250,\qquad
 \lambda^6/15120=5672362500.
\]
For every integer $m\ge1$ the following is divisibility in
$\mathbb Z[x]$:
\begin{equation}\label{eq:inverse-third}
 \mathcal T\bigl(2\psi_{1+2\lambda^2x}(m),x\bigr)-2m
 \in x^4\mathbb Z[x].
\end{equation}
Indeed, the first four recurrence coefficients are
\[
 2m+8\lambda^2\binom{m+1}{3}x
 +32\lambda^4\binom{m+2}{5}x^2
 +128\lambda^6\binom{m+3}{7}x^3.
\]
Substitution makes the coefficients of $x,x^2,x^3$ vanish as exact
polynomials in $\mathbb Q[m]$. Both full polynomials have integer
coefficients for each integer $m$, proving \eqref{eq:inverse-third}.
The supplied finite certificate checks these polynomial identities; it
does not replace the all-index recurrence proof by sampled values.

The transport argument needs a sign bound valid before index recovery.
For $k\ge2$ and $x\ge1$,
\begin{equation}\label{eq:inverse-third-sign}
 \mathcal T(k,x)\le k.
\end{equation}
To verify it, set $t=k^2\ge4$ and $v=\lambda^2x\ge1$. The expression
subtracted from $k$ is $k(k^2-4)$ times
\[
 \frac v6-\frac{v^2(27t+8)}{360}
 +\frac{v^3(675t^2-324t+160)}{15120}.
\]
The last two terms have nonnegative sum because
\[
 (675t^2-324t+160)-42(27t+8)
 =675(t-4)^2+3942(t-4)+4792>0.
\]

\subsection{Soundness in the required order}
The complete proof is retained in
\texttt{compact11/PROOF\_constant\_dilation.md}; the key dependencies are given here.
At $P_0=0$, write $z=Z/M$. Prime valuations in
\[
 (Q-1)z^2=G(QV^2-1)
\]
give positive integral roots $E=\tau^2$, $F_0=f_0^2$ and
$G=g^2$ with $g=\tau f_0$. This uses $Q=4\Delta^2U^2G$ and
$\gcd(E,F_0)=1$. Adjacent-coefficient Pell descent then gives
$y=|z|/g\in\mathbb Z_{>0}$, so $z=\pm gy$ is integral.
The explicit strong pair has
\[
 I=2Ug,\qquad f^2-\Delta I^2=1,\qquad Q=(\Delta I)^2,\qquad c^2\mid I.
\]
The retained signed all-index folding lemma yields
$c=\psi_A(p)$ for an odd $p\ge3$. The first norm yields
$k=2\psi_{1+2v}(m)$. Its index is at least two: $m=1$ would give
$k=2$ and $c=2Y+1$ strictly between the first two positive main Pell
ordinates. Since $1+2v>A$ and $c>k$, one also has $m<p$.

The free chart no longer has the parent's automatic bound $A>q^2$.
A separate argument supplies the needed replacement $c>2q^4$.
If $p\ge5$, then $c\ge16A^4-12A^2+1>4A^4>4q^4$.
If $p=3$, then $m=2$ and
\[
 k=8\lambda^2XY^2+4,\qquad c=4A^2-1,\qquad
 \frac ck>\frac{(X+1)^2}{2X+1}>\frac{X+1}{2}.
\]
Together with $c/k<Y+1$, integrality gives $2Y\ge X$, whence
$c>\lambda^2X^4-1>2q^4$.

Before the sign of $R$ is known, the uniform auxiliary estimate gives
\[
 L=bV>0,\qquad gy<3L/8,\qquad |b|/q^7<L/2.
\]
For completeness, $b=0$ would imply
$0<y^2-c^2=(c^2-1)/(Q-1)<1$. The auxiliary Pell index is odd and at
least three, so $|V|\ge4Q-3$. With $c_3=4\Delta+3\le c$ and $f>2c$,
\[
 |b|\ge4\Delta f-\frac{c+c_3}{f}>4\Delta f-1>4g,
\]
while $y<3|V|/2$. These prove the estimates.

Use the proof-only quantities
\[
 \sigma=Y/q^3,\quad \rho=e_*/q^7,\quad W=1+hL,\quad t=B/h.
\]
The carrier identity is
\[
 z=n+\kappa+\sigma-Wt-\rho.
\]
By \eqref{eq:inverse-third-sign}, $e_*\le k-R-1<b$.
If $B\le-1$, the identity gives
$n<L/2-L+|b|/q^7<0$. Thus $B\ge0$ and $C>0$.
The retained mask estimate now gives $3q<R<q^4-2$.
Hence $c>2R$, and the normalized strong-rank and signed auxiliary-index
lemmas apply to $V\equiv-c\pmod f$ and $V\equiv-R\pmod c$.
They give $p=R$.

Integral $z$ implies $q^7\mid h(e_*-q^4Y)$. Since $\gcd(q,h)=1$,
$q^4\mid e_*$. Equation \eqref{eq:inverse-third}, together with $q\mid v_0$,
gives $q^4\mid e_*-(2m-R-1)$. But $1\le m<R<q^4$, and therefore
\[
 2m=R+1.
\]
No coprimality between $q$ and $H$ has been used.

Write $R=2r+1$, so $r\ge48$. The inherited ratio theorem is uniform for
every $\lambda\ge1$, $Y\ge1$:
\[
 \frac{(X+1)^{2r}}{X^r}
 <\frac{\psi_A(2r+1)}{\psi_{1+2v}(r+1)}
 <\frac{(X+1)^{2r}}{X^r}
      \left(1+\frac1{Y(X+1)}\right)^{2r}.
\]
It implies $Y>X^r/3$ and $H>X^{r+1}$.
The independent equation $N=0$ gives $H\mid K$; the universal Pell
projector then gives
\[
 H\mid (X-2^R)(2^RX+1).
\]
The entire product has absolute value below
$4\cdot4^rX^2+8\cdot16^rX<X^{r+1}<H$, so it vanishes and $X=2^R$.

Since $q\mid X$, write $q=2^a$ with $a\ge5$. The inequality
$R>3q\ge7a$ gives $q^7\mid X$, hence $q^7\mid v_0$.
Equation \eqref{eq:inverse-third} now gives $q^7\mid e_*$.
Thus $\rho$ is integral, and the carrier gives
$Y/q^3-WB/h\in\mathbb Z$. Coprimality and $W\equiv1\pmod h$ yield
\[
 q^3\mid Y,\qquad h\mid B,\qquad B/h\ge0.
\]
The retained fixed-minus parity, half-binomial rounding and fresh native
reconstruction apply. At the entire supplied fixed tuple they produce a
finite syntactic ZFC refutation.

\subsection{Positive completeness and exact size}
Conversely, a refutation supplies the retained canonical outer data,
including $X=2^R$, $R\equiv3\pmod4$, $q^3\mid Y$ and $B/h\ge0$.
Form the new bases with $\lambda=210$ and set
\[
 c=\psi_A(R),\quad k=2\psi_{1+2v}((R+1)/2),\quad
 \tau=\chi_{1+2v}((R+1)/2).
\]
The uniform ratio theorem gives positive integers
$\eta=c-kY$ and $\zeta=k-\eta$. Projection gives positive integral
$\kappa=K/H$, while \eqref{eq:inverse-third} gives integral $\rho$.

A positive period of the determinant-one Pell matrix modulo $c^2\tau$
provides arbitrarily large pairs $(f_0,I_0)$ with $c^2\tau\mid I_0$.
Take $i=I_0/(c^2\tau)$, so $g=\tau f_0$, and choose the period multiple
large enough that $g>\kappa+\sigma+|\rho|$.
The doubled pair gives $f=2f_0^2-1$, $I=2f_0I_0$ and $S=\Delta I$.
Put
\[
 V=\chi_S(R)/S,\qquad y=\psi_S(R),\qquad
 T=\frac{V+c+Rf}{cf}.
\]
The odd quotient identities at $R\equiv3\pmod4$ make $T$ a positive
integer and give $V=f(cT-R)-c$. Finally
\[
 n=gy+(1+hL)(B/h)+\rho-\kappa-\sigma>0.
\]
These witnesses satisfy both $P_0=0$ and $N=0$.
The cofinal choice of the strong period is essential because the new
inverse correction may be negative and is larger than the old correction.

The complete highest homogeneous part, with $q_0=\beta J$ and
$k_0=\eta+\zeta$, is
\[
 [P]_{174}=
 -48\lambda^{36}q_0^{47}w^{25}Y^{56}i^{10}k_0^{32}(q_0-F)^2
 \bigl(3k_0Y-2q_0w-4\lambda\kappa Y\bigr).
\]
In particular,
\[
 [J^{49}w^{25}Y^{57}i^{10}\eta^{32}\kappa]P
 =192\lambda^{37}\beta^{49}>0.
\]
This proves exact degree $174$ at the actual fixed coefficients.
The independent checker verifies all $104$ live gates by exact polynomial
cuts, the inverse identity in $\mathbb Q[m,x]/(x^4)$, the complete
$300$-monomial leading form, all four reconstructed coefficient payloads,
and rejection of $208$ operator mutations. These finite checks support
the emitted artifacts; the preceding unbounded argument supplies the
zero-existence equivalence.

Deleting $\kappa$ and restoring $K/H$ inside the shared carrier would
suggest a ten-variable degree-$176$ formula. That shortcut is rejected:
fixed $\lambda$ and free $Y$ do not guarantee $\gcd(q,H)=1$, so its first-index
recovery is not proved.

